\documentclass[border=12pt]{standalone}
\usepackage[utf8]{inputenc}
\usepackage[T1]{fontenc}
\usepackage[french]{babel}
\usepackage[sfdefault]{noto}
\usepackage{tikz}
\usetikzlibrary{arrows.meta,positioning,calc}

\definecolor{marine}{HTML}{1F3864}
\definecolor{bleuclair}{HTML}{2E75B6}
\definecolor{vert}{HTML}{3C9A5F}
\definecolor{violet}{HTML}{7B4FA8}
\definecolor{gristexte}{HTML}{4A4A4A}
\definecolor{bordure}{HTML}{C9D3E0}
\definecolor{fondbleu}{HTML}{EAF2FB}
\definecolor{fondviolet}{HTML}{F3EFF9}
\definecolor{fondvert}{HTML}{EBF6EF}

\begin{document}
\begin{tikzpicture}[
  font=\sffamily,
  couche/.style={rounded corners=6pt, line width=1.2pt, text width=5.0cm,
                 minimum height=5.6cm, anchor=north west},
  fleche/.style={-{Latex[length=4mm,width=3mm]}, line width=1.6pt, marine},
  etiq/.style={rounded corners=3pt, inner sep=4pt, font=\sffamily\bfseries\scriptsize},
]

% ============================== EN-TETE ====================================
\node[anchor=west, font=\sffamily\bfseries\Large, text=marine] (titre) at (0,0)
      {LES TROIS COUCHES D'AUTOMATISATION};
\node[anchor=north west, font=\sffamily\small, text=gristexte] at ($(titre.south west)+(0,-0.06)$)
      {Qui fait quoi, et ce que l'on peut rejouer};

% =============================== COUCHE 1 ==================================
\node[couche, draw=bleuclair, fill=fondbleu] (c1) at (0,-1.5) {};
\node[anchor=north, font=\sffamily\bfseries\large, text=bleuclair] at ($(c1.north)+(0,-0.35)$)
     {1. SOCLE};
\node[anchor=north, font=\sffamily\bfseries\Large, text=marine] at ($(c1.north)+(0,-1.0)$)
     {Terraform};
\node[anchor=north, font=\sffamily\footnotesize, text=gristexte, align=center, text width=4.4cm]
     at ($(c1.north)+(0,-1.75)$)
     {Crée ce qui n'existe pas encore};
\draw[bordure, line width=0.8pt] ($(c1.north west)+(0.4,-2.5)$) -- ($(c1.north east)+(-0.4,-2.5)$);
\node[anchor=north west, font=\sffamily\scriptsize, text=gristexte, align=left, text width=4.3cm]
     at ($(c1.north west)+(0.4,-2.75)$)
     {\textbullet\ les 6 machines virtuelles\\[2pt]
      \textbullet\ les 7 réseaux VLAN\\[2pt]
      \textbullet\ les disques\\[2pt]
      \textbullet\ l'inventaire Ansible};
\node[etiq, fill=vert!15, text=vert, anchor=south] at ($(c1.south)+(0,0.28)$) {REJOUABLE};

% =============================== COUCHE 2 ==================================
\node[couche, draw=violet, fill=fondviolet] (c2) at (5.9,-1.5) {};
\node[anchor=north, font=\sffamily\bfseries\large, text=violet] at ($(c2.north)+(0,-0.35)$)
     {2. AMORÇAGE};
\node[anchor=north, font=\sffamily\bfseries\Large, text=marine] at ($(c2.north)+(0,-1.0)$)
     {cloud-init};
\node[anchor=north, font=\sffamily\footnotesize, text=gristexte, align=center, text width=4.4cm]
     at ($(c2.north)+(0,-1.75)$)
     {Rend la machine joignable,\\au premier démarrage};
\draw[bordure, line width=0.8pt] ($(c2.north west)+(0.4,-2.5)$) -- ($(c2.north east)+(-0.4,-2.5)$);
\node[anchor=north west, font=\sffamily\scriptsize, text=gristexte, align=left, text width=4.3cm]
     at ($(c2.north west)+(0.4,-2.75)$)
     {\textbullet\ compte et clé SSH\\[2pt]
      \textbullet\ adressage fixe du VLAN\\[2pt]
      \textbullet\ pare-feu local, SSH seul\\[2pt]
      \textbullet\ heure, swap, python3};
\node[etiq, fill=red!12, text=red!60!black, anchor=south] at ($(c2.south)+(0,0.28)$)
     {UNE SEULE FOIS};

% =============================== COUCHE 3 ==================================
\node[couche, draw=vert, fill=fondvert] (c3) at (11.8,-1.5) {};
\node[anchor=north, font=\sffamily\bfseries\large, text=vert] at ($(c3.north)+(0,-0.35)$)
     {3. CONFIGURATION};
\node[anchor=north, font=\sffamily\bfseries\Large, text=marine] at ($(c3.north)+(0,-1.0)$)
     {Ansible};
\node[anchor=north, font=\sffamily\footnotesize, text=gristexte, align=center, text width=4.4cm]
     at ($(c3.north)+(0,-1.75)$)
     {Installe et maintient dans l'état};
\draw[bordure, line width=0.8pt] ($(c3.north west)+(0.4,-2.5)$) -- ($(c3.north east)+(-0.4,-2.5)$);
\node[anchor=north west, font=\sffamily\scriptsize, text=gristexte, align=left, text width=4.3cm]
     at ($(c3.north west)+(0.4,-2.75)$)
     {\textbullet\ les 3 services métiers\\[2pt]
      \textbullet\ la supervision Centreon\\[2pt]
      \textbullet\ les 2 pare-feux\\[2pt]
      \textbullet\ le durcissement, réaffirmé};
\node[etiq, fill=vert!15, text=vert, anchor=south] at ($(c3.south)+(0,0.28)$) {REJOUABLE};

% =============================== FLECHES ===================================
\draw[fleche] ($(c1.east)+(0,-2.8)$) -- ($(c2.west)+(0,-2.8)$);
\draw[fleche] ($(c2.east)+(0,-2.8)$) -- ($(c3.west)+(0,-2.8)$);

% ============================== BAS DE PAGE ================================
\node[draw=bordure, rounded corners=5pt, line width=1pt, fill=white,
      text width=16.2cm, minimum height=1.15cm, anchor=north west] (note) at (0,-7.5) {};
\node[anchor=west, font=\sffamily\footnotesize, text=gristexte, text width=15.4cm, align=left]
     at ($(note.west)+(0.4,0)$)
     {\textbf{\color{marine}Pourquoi trois couches et pas deux.}
      Cloud-init ne s'exécute qu'au premier démarrage : il ne sait ni se rejouer, ni rendre compte,
      ni corriger une dérive. Tout ce qui doit rester vrai dans le temps appartient donc à Ansible.};

\end{tikzpicture}
\end{document}
